% Section 2. The ladder. Short rewrite of 2 Oct 2026, gated through mimesis voice=research,
% revised after Patrick's read (plain lead-in before each statistic; repository detail moved to
% the Open Practices Statement). Full specification: Supplement S7 and LADDER_SPEC.md.
% Level 2 reports three verdicts (kept, not kept, unresolved); the five regions name direction.

\section{The Validity Ladder}
\label{sec:ladder}

The ladder applies to any simulated sample in which a model answers a fixed instrument while conditioned on a persona.
One draw is one call to the model and the answer vector it returns.
A cell is one persona under one model and one prompt framing, and it holds $k$ draws whose mean score is the persona mean.
Figure~\ref{fig:ladder} summarises the gate and the four levels, with the question each asks, the statistic it reads and the use a pass supports.

The gate tests the base premise of any simulated sample, whether the model gives a persona a score consistent enough to read.
Each level above the gate supports a broader claim, moving from a single answer at level 1 to gaps between groups at level 2, the population's level at level 3 and the structure of the instrument at level 4.
A sample earns each claim separately, and the researcher reads off the claims their sample supports.
Without these checks, answers that look coherent are assumed coherent and a sample that looks typical is assumed free of bias.

A researcher deploys one model, so the ladder gives each model its own verdicts.
Pooling several models into one verdict can hide errors that run in opposite directions, as the sex gaps in Section~\ref{sec:results} show.
A pass means the test found no failure at the precision the data allow, and a level whose population reference is too imprecise to support any verdict records a stop in place of a pass or a fail.
The gate fixes $k$ for levels 2 and 3, and a failure at level 1 rules out single-draw uses only, so the other levels still run.
The rules and thresholds were fixed before any of the datasets in Section~\ref{sec:external} were analysed.
The thresholds are defaults that a user may change, and Supplement S7 gives the full specification, including estimators, bootstrap designs and multiplicity control.

\begin{figure}[tp]
% Figure 1 title (APA: number and title above the figure). The reading guide is fig1_ladder_note.tex.
% General map of the ladder, drawn by scripts/84d_fig1_ladder.py; no corpus or survey named.
\caption{The Validity Ladder}
\label{fig:ladder}

{\centering\includegraphics[width=\linewidth,height=0.71\textheight,keepaspectratio]{fig1_ladder.pdf}\par}
% Figure 1 note (below the figure).
\figurenote{Read from the gate at the bottom to level 4 at the top.
Each row gives the question a level asks, the unit it reads, its statistic and pass rule, what it is compared with, and the use a pass supports.
The gate sets $k$, the number of draws averaged into each persona mean, for levels 2 and 3.
A level-1 failure rules out single-draw uses only, so levels 2 to 4 run whatever level 1 finds.
Transgender and multiracial personas have no NHANES group of their own and never enter a group comparison.}

\end{figure}

\subsection{Gate: Regeneration Stability}
\label{sec:gate-method}

The gate asks how many draws must be averaged before a persona's score is precise enough to read.
Repeated draws of one persona differ because the model samples its answer, so each draw is the persona's own score plus draw-to-draw noise.
A generalizability study \citep{cronbach1972dependability,brennan2001gtheory} measures the two parts as variances, $\sigma^2_p$ for the spread of scores between personas and $\sigma^2_r$ for the spread of draws within one persona.
Averaging $k$ draws shrinks the noise in a persona mean, and its standard error and dependability are
\begin{equation}
\mathrm{SE}(k) = \sqrt{\sigma^2_r / k}, \qquad \phi(k) = \frac{\sigma^2_p}{\sigma^2_p + \sigma^2_r / k}.
\label{eq:phi}
\end{equation}
The standard error is the typical distance between a persona mean and the score the model would give over unlimited draws.
Dependability, from 0 to 1, is the share of the variation among persona means that reflects real differences between personas.

Persona means of $k$ draws may be read when $\mathrm{SE}(k)$ is at most 0.25 reference standard deviations in each framing, and the recommended $k$ brings it to 0.125 standard deviations.
On the PHQ-8 these tolerances are 0.98 and 0.49 points.
A single draw may be read as the persona's score only when the draw-to-draw standard deviation is itself at most 0.25 reference standard deviations.
Dependability is reported beside every verdict and required ($\phi \ge .80$) only for uses that compare or rank individual personas, because a faithful simulator inherits the population's small differences between cells.
Drawn through the persona design, real NHANES respondents reach a dependability of .70 at 30 draws. % R: analysis_brm/intermediate_panel/88_panel_verdicts.csv
A pass supports reading the persona mean of $k$ draws as the model's answer for that persona, and it says nothing about whether the answer is correct.

\subsection{Level 1: Individual Coherence}
\label{sec:l1-method}

Level 1 asks whether a single draw looks like an answer sheet a real person could give, and whether the draws are as varied as real people.
Real respondents with the same total score tend to endorse the common symptoms before the rare ones, and a minority answer in unusual patterns.
Person fit measures how typical an answer pattern is for its total.
A graded response model \citep{samejima1969graded} is fitted to the population reference with survey weights, and every respondent and every draw receives the standardised person-fit statistic $l_z^*$ \citep{drasgow1985appropriateness,snijders2001personfit,sinharay2016personfit}, which is low for an unusual pattern and high for an unusually regular one.
Each draw is then compared with real respondents who have the same total score.

The \emph{misfit ratio} is the share of a model's draws more unusual than 95\% of those respondents, divided by the 5\% expected, and the \emph{overfit ratio} is the share more regular than 95\% of them, divided by 5\%.
Both ratios equal 1 in a real population, and a ratio of 0.2 means the model gives a fifth as many such patterns as people do.
A model passes in a framing when the 90\% intervals of both ratios lie between $1/\tau$ and $\tau$.
The tolerance $\tau$ is the smallest of 1.5, 2 and 3 that covers how far every real subgroup by sex, race and ethnicity, income and survey cycle clearly departs from 1 when scored against the whole reference, with a floor of 1.5 set by the positive control (Section~\ref{sec:controls}).
On NHANES, $\tau = 1.5$. % R: l1_personfit_tolerance.csv
A model with too few patterns in both tails is called \emph{compressed}.
A pass supports using single draws as individual cases, such as exemplar vignettes or test items for a screener, and a compressed failure rules out uses that depend on the spread of real presentations.

\subsection{Level 2: Subgroup Fidelity}
\label{sec:l2-method}

Level 2 asks whether the gap between two groups, such as low- and high-income adults, has the population's size and direction.
The simulated gap $g$ compares personas that differ on one attribute and match on the rest, such as two personas identical except for income.
The population gap $\gamma$ is computed the same way in the survey, among people who match on the other attributes, and the crude survey gap is reported beside it.
The statistic is the ratio $\rho = g/\gamma$.
A ratio of 1 reproduces the population gap, 2 doubles it, 0 erases it, and a negative ratio reverses it.
Its interval is a Fieller interval, the standard interval for a ratio of two estimates \citep{fieller1954interval}, widened for the number of contrasts tested on each model, which gives 98.6\% intervals for seven contrasts.

A contrast is \vd{kept} when its whole interval lies between 0.75 and 1.25, \vd{not kept} when the interval lies wholly outside that range, and \vd{unresolved} when it overlaps it.
For a contrast that is not kept, the ratio gives the direction of the error, with ratios above 1.25 steepened, ratios from 0.25 to 0.75 attenuated, ratios near zero missing and negative ratios reversed.
Two stops are checked before any verdict.
A contrast is not read when the survey cannot tell its population gap from zero, and it is not read when the survey measures the gap too imprecisely for any simulated gap to be certified kept, unless the contrast is already clearly not kept.
A model passes level 2 when every contrast read is kept, fails when any is not kept, and is otherwise unresolved.
A pass supports reading the size and direction of group differences from the simulated sample, and it says nothing about either group's level.

\subsection{Level 3: Population Calibration}
\label{sec:l3-method}

Level 3 asks whether the simulated sample has a group's level, such as the mean PHQ-8 score of women or of low-income adults.
A persona grid is designed and not sampled, so before the comparison the personas are reweighted to the real mix of people in each group.
Each persona cell $c$ in a group $G$ is weighted by its share $p_c$ of the group in the reference, and the residual is
\begin{equation}
R_G = \sum_{c \in G} p_c\,(s_c - y_c),
\end{equation}
where $s_c$ is the persona mean and $y_c$ the mean of real respondents who match that persona \citep{holtsmith1979poststrat}.
A residual of $+2$ means the simulated group scores 2 points higher than the real group.
Its interval combines the draw error within cells with a design-based jackknife of the survey.

A row passes at tolerance $\delta$ when its 90\% interval lies inside $\pm\delta$, an equivalence test by two one-sided tests \citep{schuirmann1987tost,lakens2018tutorial}.
The researcher names $\delta$ as the smallest difference that matters for the use.
By default it is half a reference standard deviation, about 2 points on the PHQ-8, and 0.2 standard deviations is the strict alternative.
A model passes when every group row passes.
The share scoring at or above the instrument's cut point is a second row, because a correct mean can hide a wrong prevalence.
A pass supports using the simulated level as a stand-in for the group's level, within $\delta$, for the prompt, persona grid and survey period tested.

\subsection{Level 4: Structural Fidelity}
\label{sec:l4-method}

Level 4 asks whether the items of the instrument hang together as they do in people.
In real respondents the eight PHQ-8 items rise and fall together, so the total works as one depression score, and it works the same way in different groups.
All draws of one model in one framing are treated as one population, with personas as the resampling unit.
A one-factor model is fitted to the item correlations \citep{olsson1979polychoric} and checked against three rules.
R1 requires a general factor, one shared dimension behind the items, with every loading at least .30 and the first eigenvalue at least three times the second, or the reference's own value where the reference falls short.
R2 requires the items to load on that factor in the same pattern and with the same strength as in people, read as Tucker's congruence of at least .95 \citep{lorenzoseva2006congruence} and a root mean square difference in loadings of at most .10, each at the conservative end of its 90\% interval.
R4 requires the factor to work the same way across sex, race and income wherever it does in people, read on the RMSEA of each invariance step \citep{savalei2024rmsea} against a margin of .08.

A model passes level 4 when it passes all three rules in every framing.
Passing R1 and R2 supports scoring the total as one score with item weights like people's.
Passing R4 adds that a group comparison of the total means the same in the simulation as in people, so passes at levels 2 and 3 are read beside it.
A diagnostic (R3) also reports whether the factor holds among draws of one persona, as it does among people who share a demographic profile.
